\section{Methodology}
\label{sec:methodology}

\subsection{ConvertAsset Pipeline}
\label{sec:pipeline}

ConvertAsset is an open-source toolkit that converts NVIDIA Isaac Sim assets
from MDL materials to UsdPreviewSurface, preserving the compositional structure
of USD scenes. The pipeline operates in three stages.

\paragraph{Stage 1: Recursive USD Traversal.}
The processor begins at a root USD file and recursively discovers all
referenced assets --- sublayers, references, payloads, variant selections, and
value clips --- building a dependency graph.  A cycle-detection mechanism
(\texttt{in\_stack} set) prevents infinite recursion on circular references,
while a deduplication dictionary (\texttt{done} map) ensures each asset file is
processed exactly once. Crucially, the pipeline does \emph{not} flatten the USD
composition: each file is processed independently and saved as a sibling
\texttt{*\_noMDL.usd} file, so the original composition arcs (references,
payloads, variants) remain intact with only their asset paths rewritten to
point to the converted counterparts.

\paragraph{Stage 2: Material Conversion.}
For each USD file in the dependency graph, the pipeline identifies all
materials driven by MDL shaders. It constructs a minimal UsdPreviewSurface
shader network per material, mapping four primary channels: base color (diffuse
albedo), roughness, metallic, and normal. Texture assets referenced by the MDL
shader are located, resolved against the declaring layer's directory, and
connected to the new UsdPreviewSurface inputs via \texttt{UsdShadeConnectableAPI}.
When texture files cannot be resolved, constant fallback values are used.
MDL features not captured by the four-channel mapping include subsurface
scattering, clearcoat layers, anisotropic reflectance, emission, opacity, and
displacement. After the UsdPreviewSurface network is established, all MDL shader prims and
MDL-specific material outputs are removed. A post-conversion verification pass
confirms that no MDL residuals remain.

\paragraph{Stage 3: Asset Path Rewriting.}
All composition arcs in the converted file --- references, payloads, sublayer
paths, and variant-internal asset paths --- are rewritten from their original
targets to the corresponding \texttt{*\_noMDL} variants. This ensures that
opening the converted root file yields a fully self-consistent scene composed
entirely of UsdPreviewSurface materials.

\subsection{Experimental Setup}
\label{sec:setup}

\paragraph{Dataset.}
We select four chest-of-drawers assets from the NVIDIA Isaac Sim asset library:
\texttt{chestofdrawers\_0004}, \texttt{chestofdrawers\_0011},
\texttt{chestofdrawers\_0023}, and \texttt{chestofdrawers\_0029}. These objects
feature diverse wood, metal, and fabric materials with varying degrees of MDL
complexity, including multi-layer coatings and procedural textures, making them
representative test cases for material simplification.

\paragraph{Rendering.}
For each object, we render images from six camera viewpoints distributed around
the scene: front, back, left, right, top-front-left, and top-front-right. Each viewpoint
is rendered twice --- once with the original MDL materials and once with the
ConvertAsset-converted UsdPreviewSurface materials --- yielding $4 \times 6
\times 2 = 48$ images total (24 matched pairs). All renders use identical
geometry, lighting, and camera parameters at $1024 \times 1024$ resolution to
isolate the effect of material representation.

\paragraph{Hardware.}
All experiments are conducted on a workstation equipped with an NVIDIA GeForce
RTX 4090 (24\,GB VRAM) GPU and an Intel Xeon Gold 6462C processor. Rendering uses NVIDIA Isaac Sim's
built-in RTX renderer. Feature extraction and detection inference run on the
same GPU hardware.

\subsection{Evaluation Metrics}
\label{sec:metrics}

We evaluate the impact of material simplification at four complementary levels.

\paragraph{Rendering Performance.}
We measure scene load time (seconds from \texttt{open\_stage} to first rendered
frame), peak GPU memory consumption (MB), and rendering throughput (frames per
second) for both MDL and UsdPreviewSurface versions of each scene.

\paragraph{Pixel-Level Image Quality.}
For each of the 24 matched image pairs, we compute:
\begin{itemize}
    \item \textbf{PSNR} (Peak Signal-to-Noise Ratio): measures pixel-level
    reconstruction accuracy in decibels (higher is better).
    \item \textbf{SSIM}~\cite{Wang2004Image} (Structural Similarity Index):
    captures luminance, contrast, and structural distortion on a $[0,1]$ scale
    (higher is better).
    \item \textbf{LPIPS}~\cite{Zhang2018Unreasonable} (Learned Perceptual
    Image Patch Similarity, VGG backbone): uses deep features to measure
    perceptual distance (lower is better).
\end{itemize}

\paragraph{Feature-Level Preservation.}
To assess whether material simplification alters semantic content as perceived
by modern vision models, we extract embeddings from both
CLIP~\cite{Radford2021Learning} (ViT-B/32) and
DINOv2~\cite{Oquab2023DINOv2} (ViT-B/14) for each matched pair and compute
cosine similarity. High similarity ($> 0.95$) indicates that the feature
representations are largely invariant to the material change. We report
per-pair scores and across-scene averages to identify whether specific material
types are more sensitive to simplification.

\paragraph{Detection Proxy.}
As a downstream task probe, we apply a pretrained YOLOv8 object detector
(nano variant, YOLOv8n), selected as a representative lightweight detector
commonly deployed in robotics applications, in a
zero-shot setting (no fine-tuning on our rendered data) to both MDL and
UsdPreviewSurface images. We compare detection confidence scores for the
target object across matched pairs. A statistically significant drop in
confidence would indicate that material simplification degrades object
appearance sufficiently to affect recognition, even for a detector trained on
real-world images.
